\section{Summary}
\label{sec:eval-summary}

The evidence in this chapter addresses the four research questions as follows.
For modularity (RQ1) the dependency boundary is build-enforced and checked by an executable test, an in-process module is added in one registration, and three external providers connect with no change to the core, one of them wrapping a third party's model.
For the sensing pipeline (RQ2) the eye tracker streamed at its rated \SI{90}{\hertz} with validity above 92\%, and the client transport latency stayed within the \SI{100}{\milli\second} budget for every sample, while the automated decision-path latencies remained instrumented but unexercised.
For context preservation (RQ3) the mechanism ran and recorded an audited outcome for every intervention, the reading-resume time and the post-intervention regression rate were both lower with preservation than without, and the restore over-repositioned on small reflows.
For researcher control (RQ4) the sessions were operated through a single gated console, and the complete results of this chapter were reconstructed from the exported records alone.
\Cref{ch:discussion} interprets these results, compares them with the problem statement and prior work, and develops the limitations they expose.
